\begin{proposition}[C.3.13 (Proved): a corrected marked Seifert comparison]
There is a smooth map $\mathcal G:X_c\to S^3$, invariant under
the $\widehat\delta$ circle, with the following properties.
It covers the given base coordinate through the explicit weighted
Seifert map
\[
 p:S^3\longrightarrow\PP^1,\qquad
 p(z_1,z_2)=\frac{z_2^3}{z_2^3-z_1^4},\qquad
 \lambda\cdot(z_1,z_2)=(\lambda^3z_1,\lambda^4z_2).
\]
Outside a chosen small cusp body it is the actual quotient by
the regular $\ker\gamma$ subtori, followed by a specified
base-dependent circle gauge and the explicit identification below.
It includes the actual finite root quotients with rotation data
$+1/3$ and $-1/4$. Inside the cusp body the map is modified;
it does not preserve the original translational leaves.

On the outer cusp collar a positively marked $\widehat\gamma$
circle maps orientation-preservingly onto a regular weighted
Seifert fibre. The central fibre of the target is the positive
$(3,4)$ torus knot, and a nearby regular fibre has linking number
$12$ with it, in the standard orientation of $S^3\subset\C^2$.
This is framing data in the constructed three-dimensional target.
No identification of its normal framing with the six-dimensional
orbit-knot framing, and no value of $6[\Phi]$, follows from this
statement.
\end{proposition}
\begin{proof}
We first fix the real character and all of its overlap corrections.
On an ordinary fibre, let
$\gamma_z:\C^2\to\R$ be the real-linear functional determined by
\[
 \gamma_z(\Pi_c(z)v)=\gamma(v)\qquad(v\in\R^4).
\]
The real period matrix is invertible in the admissible range.
The resulting circle coordinate
$\theta=\gamma_z(\zeta)\pmod\Z$ is invariant under ordinary
monodromy because $\gamma A_j=\gamma$. It is the literal
regular-fibre quotient character, not a new choice of marking.

On the $j$th finite root chart write $\theta_j$ for this
character before the affine finite quotient. The exact map from
that chart to the ordinary chart is translation by
$\sigma_{j,c}=w_j\Pi_cv_j$, where
$w_j=(2\pi i)^{-1}\log s_j$ and
$\epsilon_j=\gamma(v_j)$ is $1,-1$, respectively.
Consequently its real circle overlap is exactly
\begin{equation}\label{eq:C13-real-overlap}
 \theta=\theta_j+\epsilon_j\operatorname{Re}w_j+a_j^{\mathrm{rad}}(t),
 \qquad
 a_j^{\mathrm{rad}}(t)=\operatorname{Im}w_j\,
                  \gamma_z(i\Pi_cv_j).
\end{equation}
The last expression is real and single-valued on the punctured
base disc. Indeed $\operatorname{Im}w_j$ is single-valued,
$v_j$ is fixed by the local integral monodromy, and its
complex-linear covariance preserves $\gamma_z$.
No logarithmic radial term has been dropped.

Choose disjoint finite discs and smooth radial cutoffs equal to
one on smaller discs and zero outside the larger ones. They define
a real smooth function $d(t)$ on $B^\circ$ which is $-a_j^{\mathrm{rad}}(t)$
on the smaller punctured discs and is zero on a cusp collar.
Choose the roots $F^3=t$, $G^4=t-1$ on the uniformizing cover and
the logarithm $\ell=(2\pi i)^{-1}\log(F/G)$ used in C.3.5.
Put
\[
 \Theta=\theta-\operatorname{Re}\ell+d(t).
\]
Let $R_t>0$ be the unique solution of
\[
 R_t^6|G|^2+R_t^8|F|^2=1.
\]
Existence and uniqueness follow from strict increase from zero
to infinity; the derivative in $R_t$ is positive, so the
implicit-function theorem gives the required smooth dependence.
The proposed map on the ordinary region is
\begin{equation}\label{eq:C13-weighted-map}
 (z_1,z_2)=
 \bigl(R_t^3G e^{6\pi i\Theta},\,
                R_t^4F e^{8\pi i\Theta}\bigr).
\end{equation}
Its norm is one and direct substitution gives $p(z_1,z_2)=t$.
Increasing $\theta$ by $a$ is exactly the weighted action of
$e^{2\pi ia}$, so the positive $\gamma$ orientation is retained.

We check descent and the finite centres explicitly. For the
clockwise first root generator $F$ changes by $e^{-2\pi i/3}$,
$G$ is unchanged, and $\ell$ changes by $-1/3$.
Thus $\Theta$ changes by $+1/3$, and both coordinates in
\eqref{eq:C13-weighted-map} remain unchanged. For the second
generator $G$ changes by $e^{-2\pi i/4}$, $F$ is unchanged,
and $\Theta$ changes by $-1/4$. Again both coordinates are
unchanged: their total exponents are respectively $-1$ and
$-1$. These generators and changes of logarithm prove descent
on the entire ordinary base.

On a full finite root disc the exact identity from C.3.5 is
$\ell=\epsilon_jw_j+\kappa_j$, where $\kappa_j$ is a
holomorphic logarithm of an invariant nonvanishing unit.
Equations \eqref{eq:C13-real-overlap} and $d=-a_j^{\mathrm{rad}}$ therefore give
\[
 \Theta=\theta_j-\operatorname{Re}\kappa_j.
\]
This expression is smooth through the root centre. The vanishing
factor $F$ at the first centre, or $G$ at the second, is the
root coordinate times a holomorphic unit. The other factor is
nonzero. Thus \eqref{eq:C13-weighted-map} extends smoothly to
the full root chart. Its transverse derivative is nonzero and
the circle factor has degree one; it identifies the actual
finite quotient solid torus with the corresponding weighted
solid torus. The deck actions on $(s_j,\theta_j)$ are precisely
\[
 (s_1,\theta_1)\mapsto(e^{-2\pi i/3}s_1,\theta_1+1/3),
 \qquad
 (s_2,\theta_2)\mapsto(e^{-2\pi i/4}s_2,\theta_2-1/4).
\]
This verifies the overlap maps, not merely the exceptional
multiplicities of an abstract Seifert space.

We next verify the cusp meridian and construct the replacement
map upstairs, without assuming a smooth quotient structure at
the fixed set. Let $V$ be a small smooth compact cusp body and
$\overline V=V/\rho$. C.3.7--C.3.9 identify
$H_1(\overline V;\Z)=\Z\widehat\gamma\oplus\Z\widehat u$.
The outer orbit link is the mapping torus of the action
$u\mapsto u+w$ on the $(u,w)$ torus, times the $\gamma$ circle.
Its first homology is therefore freely generated by
$\widehat\gamma,\widehat u$ and a transverse cusp meridian;
$\widehat w$ is a commutator. The two first generators have
the stated images in $H_1(\overline V)$ by the actual fan marking
in C.3.9.

The transverse meridian has zero $\gamma$ value and bounds
an actual disc in $V$. To see this without choosing a meridian
up to an unknown fibre twist, take the toric section
$z_1=z_2=1$. In the chart of C.3.12 it is
$(x,y,z)=(1,q,1)$, so it extends across $q=0$.
Its additive logarithms are zero and hence $\theta=0$
identically on its boundary. Thus the integral class on
$\overline V$ with values $1,0$ on
$(\widehat\gamma,\widehat u)$ restricts to exactly the raw
character $[d\theta]$ on the orbit link, including the transverse
meridian value. Pull it back to $V$.

Represent the pulled-back integral class by a smooth closed real
one-form $\alpha$ with integral periods. Average it over the
$\delta$ circle. Each group element is isotopic to the identity,
so the cohomology class and its integral periods are unchanged.
The averaged form is invariant. If $K$ denotes the circle
generator, Cartan's formula gives $d(\alpha(K))=0$.
The body is connected and has fixed points, where $K=0$;
therefore $\alpha(K)=0$ everywhere.
On an outer collar $\alpha-d\theta$ is exact, since the two
forms have the same periods. Choose its real primitive,
average the primitive over the circle, and extend it to $V$
using an invariant collar cutoff. Subtracting its differential
from $\alpha$ makes the resulting invariant integral closed
form equal to $d\theta$ on a smaller outer collar. Integrating
it modulo $\Z$, with the constant chosen at one collar point,
gives a smooth map
\[
 g:V\longrightarrow\R/\Z
\]
which equals $\theta$ on that collar. Its contraction with $K$
is zero, so $g$ is invariant. This construction is entirely on
the smooth manifold $V$ and yields a continuous factor on
$\overline V$.

For completeness, the prescribed orbit boundary map has a
unique relative homotopy class of extensions to $\overline V$.
The restriction of $H^1(\overline V;\Z)$ to the link is
injective by the two period generators just checked, and both
spaces are connected. The relative cohomology exact sequence
therefore gives $H^1(\overline V,\partial\overline V;\Z)=0$.
Maps to $\R/\Z$ with fixed boundary values are classified by
this relative group: a map represents its integral period class,
and a map with zero class has a real lift which supplies the
relative homotopy. This assertion does not classify normal
framings of any embedded trace.

Use now the holomorphic cusp coordinate $q=1/t$; changing from
the original Mumford coordinate multiplies it by a unit and
does not change the meridian just verified. Write
$q=e^{2\pi i s}$. Compatible roots satisfy
\[
 F=q^{-1/3},\qquad G=q^{-1/4}(1-q)^{1/4},\qquad
 \ell=-s/12+\kappa_0(q),\qquad
 \kappa_0(q)=-\frac{\log(1-q)}{8\pi i}.
\]
A different compatible root choice only changes the displayed
unit gauge. The smooth fibre coordinate in the weighted model
near its regular cusp knot is
\begin{equation}\label{eq:C13-cusp-gauge}
 \vartheta=\Theta-\operatorname{Re}s/12
           =\theta-\operatorname{Re}\kappa_0(q),
\end{equation}
because $d=0$ there. In particular no additional integer
meridian winding is present. This is the required boundary
framing comparison; the finite labels alone would not give it.

Let $S(q)>0$ solve
$S(q)^6|1-q|^{1/2}+S(q)^8=1$. The explicit smooth solid-torus
chart of the target is
\[
 \Phi_0(q,\vartheta)=
 \bigl(S(q)^3(1-q)^{1/4}e^{6\pi i\vartheta},
                         S(q)^4e^{8\pi i\vartheta}\bigr).
\]
It is a diffeomorphism onto a neighbourhood of the regular
cusp knot. Indeed its inverse has
$q=1-z_1^4/z_2^3$ and
$e^{2\pi i\vartheta}=z_2(1-q)^{1/4}/(S(q)z_1)$;
all factors are nonzero there. Define the replacement on $V$
by
\[
 \mathcal G(x)=
 \Phi_0\bigl(q(x),g(x)-\operatorname{Re}\kappa_0(q(x))\bigr).
\]
It agrees exactly with \eqref{eq:C13-weighted-map} on the outer
collar by \eqref{eq:C13-cusp-gauge}, and hence glues with the
ordinary and finite maps already constructed. Smoothness and
$\delta$ invariance hold upstairs on every piece. The first
coordinate of the solid-torus chart proves that it still covers
$f$. By C.3.12 its new cusp leaves cannot be the original
translational leaves on every punctured fibre; that change is
an explicit part of the construction.

Finally the knot $q=0$ in the target is
$z_2^3=z_1^4$, parametrized with positive orientation by the
weighted circle $(e^{3i\phi},e^{4i\phi})$ with the fixed positive
radii prescribed by $S(0)$. It is the positive $(3,4)$ torus
knot. For a small fixed $q\ne0$, along the nearby fibre
\[
 (z_2^3-z_1^4)\circ\Phi_0(q,\vartheta)
                  =q S(q)^{12}e^{24\pi i\vartheta}.
\]
The argument therefore has degree $12$. The function
$z_2^3-z_1^4$ has a simple complex-normal zero along the
knot in $S^3$, so its argument represents the positive
meridian class in the knot complement. This degree is its
linking number with the knot. The orientation is the one
already matched to positive $\gamma$ in
\eqref{eq:C13-weighted-map}. The computation establishes the
three-dimensional push-off datum only; the normal correction
of C.3.11 remains a separate unresolved comparison.
The finite phases, radial cancellation, link abelianization and
polynomial identities are checked exactly in
\path{code/d_push_corrected_seifert.py}.
\end{proof}
